% glossary_entries.tex
\newglossaryentry{adas}{
    name={Advanced Driver Assistance System (ADAS)},
    description={Systèmes permettant d'assister le conducteur lors de la conduite}
}

\newglossaryentry{ros}{
    name={Robot Operating System (ROS)},
    description={Plateforme de développement logicielle Open Source contenant de nombreuses bibliothèques et outils pour des applications robotiques}
}

\newglossaryentry{compiler}{
    name={Compiler},
    description={Traduire du code source compréhensible par un humain en code machine}
}

\newglossaryentry{segmentation}{
    name={Segmentation d'image},
    description={Technique de vision artificielle visant à diviser une image en plusieurs régions. Chaque segment correspond généralement à une partie distincte de l'image, facilitant ainsi l'identification et la classification des objets ou des zones d'intérêt}
}

\newglossaryentry{brique mobilex}{
    name={Brique MOBILEX},
    description={Brique technologique assurant la fonction de gestion de la trajectoire du robot en téléopération assistée ou en conduite autonome}
}

\newglossaryentry{mobilex}{
    name={Challenge MOBILEX},
    description={Challenge organisé par la DGA visant à développer des fonctions de conduite autonome et d'aide à la conduite pour un robot navigant dans un environnement complexe}
}

\newglossaryentry{envcomp}{
    name={Environnement complexe/Environnement non structuré},
    description={Environnement n'obéissant à aucune règles particulières, pouvant contenir des obstacles de toutes natures et toutes formes}
}

\newglossaryentry{inference}{
    name={Inférence},
    description={L'inférence en Intelligence Artificielle représente le processus par lequel un modèle d'IA utilise ce qu'il a appris pendant l'entraînement pour faire des prédictions ou prendre des décisions sur de nouvelles données}
}

\newglossaryentry{kalman}{
    name={Filtre de Kalman},
    description={Le filtre de Kalman est un algorithme qui utilise une série de mesures observées au fil du temps, contenant des bruits et des incertitudes, et produit des estimations des variables inconnues qui tendent à être plus précises que celles basées sur une seule mesure}
}

\newglossaryentry{lidar1}{
    name={LIDAR},
    description={Capteur utilisant la lumière sous la forme d'un laser pulsé pour mesurer des distances de manière précise}
}

\newglossaryentry{platooning}{
    name={Platooning},
    description={Stratégie d'optimisation de flottes de véhicules consistant à faire circuler plusieurs véhicules en convoi rapproché pour réduire la consommation de carburant et la pollution}
}